\chapter*{Document History}
\addcontentsline{toc}{chapter}{Document History}
\index{document history}
\section*{How to Use This Chapter}
Add the newest approved change at the top of the table. Retain every approved
entry so that the document's evolution and review history remain auditable.
Every history entry shall link to the chapter, section, table, figure, or
controlled artifact where the documented change occurred.

\noindent\textbf{Author initials:} \textbf{SAA}: Stel Angelidis;
\textbf{CLAUDE}: Anthropic Claude Code; \textbf{DDM}: David Darian Muresan;
\textbf{CDX}: OpenAI Codex. A comma identifies joint authorship or implementation.

\noindent Entries from the template baseline through version 2.7 record the
evolution of the master architecture template from which \ProjectName\ was
initialized. Project history begins at version 0.1.

\begin{tabularx}{\linewidth}{L{3.1cm}L{2.2cm}YL{2.0cm}}
\caption{Document Update History\label{Table:DocumentHistory}}\\
\toprule
\textbf{Date / Version} & \textbf{Author Initials} & \textbf{Change} & \textbf{Changed At}\\
\midrule
\endfirsthead
\multicolumn{4}{l}{\small\itshape Document Update History (continued)}\\
\toprule
\textbf{Date / Version} & \textbf{Author Initials} & \textbf{Change} & \textbf{Changed At}\\
\midrule
\endhead
\midrule
\multicolumn{4}{r}{\small\itshape Continued on next page}\\
\endfoot
\bottomrule
\endlastfoot
October 14, 2026 / 0.1 & SAA, CLAUDE & Initialized \ProjectName\ from template version 2.7: applied the \texttt{tgm} file namespace to all project-owned files and TikZ sources; set project name, owner, organization, version, status, and date; recorded the Embedded Systems profile as Applicable on a microcontroller platform; recorded Stel Angelidis as authorized approver; recorded the system purpose, the decision to provide both a physical controller and a phone application, the stakeholders, the privacy, USD 20 component-cost, and microcontroller constraints, and open context items; added the IoT acronym; added a Claude Code repository bootstrap and the initialization prompt record. & \ChapterRef{Introduction}, \hyperref[Section:DevelopmentPlanRoles]{approval authority}, \hyperref[Section:DevelopmentPlanProfiles]{profile decisions}, \hyperref[Section:SystemPurpose]{system purpose}, \hyperref[Section:Stakeholders]{stakeholders}, \hyperref[Section:StakeholderConstraints]{constraints}, \ChapterRef{AIPromptLog}\\
September 18, 2026 / 2.7 & DDM, CDX & Required every applicable domain-profile concern to map to a new or existing project requirement, or to an authorized Not Applicable decision; added AI task routing and requirements-work instructions for this translation. & \hyperref[Section:DomainProfileRequirements]{Profile requirement rule}, \ChapterRef{AIAgentGuide}, \ChapterRef{EmbeddedSystemsProfile}\\
September 18, 2026 / 2.6 & DDM, CDX & Added the domain-tailoring framework and Embedded Systems Tailoring Profile; added profile selection to initialization and the Development Plan; and extended AI routing, workflows, prompts, and completion checks for applicable profiles. & \hyperref[Section:DomainProfiles]{Domain profiles}, \ChapterRef{EmbeddedSystemsProfile}, \hyperref[Section:DevelopmentPlanProfiles]{profile decisions}, \ChapterRef{AIAgentGuide}\\
September 9, 2026 / 2.5 & DDM, CDX & Required the project-initialization prompt to confirm the repository's AI tools before selecting bootstrap files and to create \texttt{SKILL.md} only for a confirmed reusable AI skill. & \hyperref[Section:InitializeNewProject]{Project initialization}\\
September 9, 2026 / 2.4 & DDM, CDX & Added an explicit copy-and-paste AI request to the project-initialization section, including namespace confirmation, evidence review, gap analysis, human review, controlled bootstrap creation, and build verification. & \hyperref[Section:InitializeNewProject]{Project initialization}\\
September 9, 2026 / 2.3 & DDM, CDX & Made TikZ the exclusive diagram source format; removed Graphviz generation from both platform build files and removed current fallback, installation, build, troubleshooting, and AI guidance for Graphviz. & \hyperref[Section:TikZAuthoring]{TikZ authoring}, \ChapterRef{DocumentGuide}, \ChapterRef{AIAgentGuide}\\
September 9, 2026 / 2.2 & DDM, CDX & Added a concise Introduction workflow for instantiating a project-specific architecture book from the master template, including namespace confirmation, evidence review, dependency-ordered development, AI bootstrap guidance, and build verification. & \hyperref[Section:InitializeNewProject]{Project initialization}, \ChapterRef{AIAgentGuide}\\
September 9, 2026 / 2.1 & DDM, CDX & Added the repository AI bootstrap procedure, applicable discovery-file catalog, non-duplication rule, reusable creation prompt, tool-mandated filename exception, and completion check. & \hyperref[Section:AIRepositoryBootstrap]{AI bootstrap}, \hyperref[Table:DocumentAuthoringRequirements]{authoring requirements}, \hyperref[Section:AIExamplePrompts]{example prompts}, \hyperref[Table:AICompletionChecklist]{completion checklist}\\
September 8, 2026 / 2.0 & DDM, CDX & Required a human-confirmed project namespace for every new project-owned file; required each TikZ directory to exactly match its including LaTeX file; and renamed both active TikZ sources to comply. & \hyperref[Table:DocumentAuthoringRequirements]{authoring requirements}, \hyperref[Section:TikZAuthoring]{TikZ authoring}, \ChapterRef{AIAgentGuide}, \hyperref[Figure:VModel]{V-model}, \hyperref[Figure:UseCaseDiagram]{use-case diagram}\\
September 7, 2026 / 1.9 & DDM, CDX & Linked each use-case identifier in the high-level TikZ diagram directly to its authoritative use-case definition. & \hyperref[Figure:UseCaseDiagram]{use-case diagram}\\
September 7, 2026 / 1.8 & DDM, CDX & Replaced the active V-model and high-level use-case Graphviz figures with chapter-owned TikZ sources and updated both chapters to include and reference the vector diagrams directly. & \hyperref[Figure:VModel]{V-model}, \hyperref[Figure:UseCaseDiagram]{use-case diagram}\\
September 7, 2026 / 1.7 & DDM, CDX & Established TikZ as the primary diagram-authoring tool, retained Graphviz as a fallback, documented the source and inclusion pattern, and added the mandatory AI rule for locating and referencing every TikZ source. & \ChapterRef{DocumentGuide}, \ChapterRef{AIAgentGuide}\\
September 6, 2026 / 1.6 & DDM, CDX & Added the requirement identifier grammar, active family-code definitions, optional project scope, stable sequence rule, classification boundary, and the controlled AI procedure for proposing and approving a new scope or family without changing existing IDs. & \hyperref[Table:RequirementIdentifierConvention]{ID convention}, \ChapterRef{AIAgentGuide}\\
August 25, 2026 / 1.5 & DDM, CDX & Expanded every formal table to the available portrait or landscape width using flexible multipage columns; retained the existing column hierarchy while increasing space for descriptive content. & \hyperref[Table:DocumentHistory]{History}, \ChapterRef{Requirements}, \ChapterRef{UserStories}, \ChapterRef{Tests}, \ChapterRef{Traceability}, \ChapterRef{Verification}, \ChapterRef{Validation}, \ChapterRef{AIAgentGuide}\\
August 25, 2026 / 1.4 & DDM, CDX & Added a dedicated AI Agent Usage Guide with efficient task routing, in-document context packs, workflows, checklists, human-review boundaries, and reusable prompts; added a concise AI entry point in the Introduction; and consolidated duplicated guidance. & \ChapterRef{AIAgentGuide}, \ChapterRef{Introduction}, \ChapterRef{DocumentGuide}\\
August 25, 2026 / 1.3 & DDM, CDX & Reorganized the requirement catalog into explicit families and subtypes; classified scalability, reliability, security, and related characteristics as quality attributes; and clarified the role of specialized working registers. & \hyperref[Table:RequirementCategoryCatalog]{Catalog}, \ChapterRef{Requirements}\\
August 25, 2026 / 1.2 & DDM, CDX & Added the Open, In Progress, Done, and Retired requirement lifecycle; defined its relationship to priority and implementation version; and added Status to every requirement register. & \hyperref[Table:RequirementStatusDefinitions]{Status}, \ChapterRef{Requirements}\\
August 25, 2026 / 1.1 & DDM, CDX & Split requirements into category-specific registers; added business requirements, a broad category catalog, functional-versus-quality guidance, and instructions for project-specific categories; and completed reciprocal mappings. & \hyperref[Table:BusinessRequirements]{B}, \hyperref[Table:FunctionalRequirements]{F}, \hyperref[Table:QualityRequirements]{Q}, \hyperref[Table:InterfaceRequirements]{I}, \hyperref[Table:SecurityRequirements]{S}, \hyperref[Table:DataRequirements]{D}, \hyperref[Table:ConstraintRequirements]{C}, \hyperref[Table:RequirementCategoryCatalog]{Catalog}\\
August 25, 2026 / 1.0 & DDM, CDX & Added explicit generic interface, security, data, and constraint requirement examples with reciprocal story, use-case, test, traceability, and verification mappings. & \ChapterRef{Requirements}, \hyperref[Table:UserStories]{US}, \hyperref[Chapter:UseCases]{UC}, \hyperref[Chapter:Tests]{TC}, \hyperref[Table:TraceabilityMatrix]{RTM}, \hyperref[Table:VerificationRegister]{VER}\\
August 25, 2026 / 0.9 & DDM, CDX & Removed all end-input boundaries and obsolete trailing project content; retained only generic template material; and added an AI-assisted architecture workflow with human-review and prompt-log guidance. & \ChapterRef{Introduction}, \hyperref[Section:AIArchitectureWorkflow]{AI workflow}, \ChapterRef{AIPromptLog}\\
August 25, 2026 / 0.8 & DDM, CDX & Changed the joint-author separator from a slash to a comma. & \hyperref[Table:DocumentHistory]{History}\\
August 25, 2026 / 0.7 & DDM, CDX & Replaced template-originator initials with named author initials and documented Codex joint authorship and implementation. & \hyperref[Table:DocumentHistory]{History}\\
August 25, 2026 / 0.6 & DDM, CDX & Added visible template attribution, author contact information, AI-agent development credit, PDF creator metadata, source-header provenance, and authoring-guide provenance guidance. & \hyperref[Chapter:TemplateAttribution]{Template attribution}, \ChapterRef{DocumentGuide}\\
August 25, 2026 / 0.5 & DDM, CDX & Replaced the literal x in the requirements implementation-version field with the empty-set symbol to identify requirements with no assigned implementation version. & \ChapterRef{Requirements}\\
August 25, 2026 / 0.4 & DDM, CDX & Expanded the generic introduction with the central lesson, source-of-truth guidance, and V-model; added a generic high-level use-case diagram; and expanded the authoring and tool-installation appendix. & \ChapterRef{Introduction}, \hyperref[Figure:VModel]{V-model}, \ChapterRef{UseCases}, \hyperref[Figure:UseCaseDiagram]{Use-case diagram}, \ChapterRef{DocumentGuide}\\
August 25, 2026 / 0.3 & DDM, CDX & Added linked history records; glossary-linked priority terms; user-story, traceability, verification, and validation registers; implementation-version status; and a final index chapter. & \hyperref[Table:DocumentHistory]{History}, \ChapterRef{Introduction}, \ChapterRef{Requirements}, \hyperref[Table:TraceabilityMatrix]{Traceability}, \hyperref[Table:VerificationRegister]{Verification}, \hyperref[Table:ValidationRegister]{Validation}, \hyperref[Chapter:Index]{Index}\\
August 24, 2026 / 0.2 & DDM, CDX & Added longtable history and requirements registers, priority and implemented-version fields, reciprocal use-case mappings, and test-case mappings. & \hyperref[Table:DocumentHistory]{History}, \ChapterRef{Requirements}, \ChapterRef{UseCases}, \ChapterRef{Tests}\\
\textit{[baseline date / 0.1]} & DDM & Generic software architecture design template baseline. & \ChapterRef{Introduction}\\
\end{tabularx}
